\subsection{All fusion rings versus categorifiable fusion rings}
There are two distinct ways an unrestricted classification could be
completed. One could parametrize all rank-four fusion rings and then
prove a uniform realization-or-exclusion theorem for those parameters.
Alternatively one could bound the multiplicity or Frobenius--Perron
dimension of every potentially categorifiable ring and use the bounded
classification. The latter route does not require classifying all
noncategorifiable fusion rings.

\begin{proposition}[An infinite family of raw fusion rings]\label{prop:raw-family}
For all integers $a,b\geq0$, the tensor product of the rank-two based rings
$x^2=1+ax$ and $y^2=1+by$ is a self-dual rank-four fusion ring. With
$z=xy$, its remaining products are
\[
 xz=y+az,\qquad yz=x+bz,\qquad
 z^2=1+ax+by+abz.
\]
Its multiplicity is $\max\{1,a,b,ab\}$. It admits a complex fusion
categorification if and only if $a,b\in\{0,1\}$.
\end{proposition}
\begin{proof}
The tensor product basis is $1,x,y,xy$; positivity, the unit coefficient,
duality and associativity are inherited from the two rank-two based rings.
The displayed products give the multiplicity. A categorification would
categorify both rank-two subrings, forcing $a,b\in\{0,1\}$ by \cite[Proposition~\ref{B-prop:ranktwo}]{Finite}. Conversely their pointed or Fibonacci unitary models
have the required Deligne products.
\end{proof}
In particular, raw rank-four fusion rings have unbounded multiplicity and
cannot be classified by a finite list. This is not an obstruction to a
parametric classification: the entire infinite family above is decided by
one rank-two theorem.

\subsection{Diophantine presentations and finite census data}
Appendix~\ref{B-sec:census} of \cite{Finite} recalls the self-dual
ten-parameter and non-self-dual six-parameter tensor presentations and
their necessary and sufficient associativity equations. The
bounded classification \cite{Finite} uses pinned complete repository
data through multiplicity $128$, comprising $889530$ rings. These
bounded lists do not parametrize the unbounded nonnegative solutions
in disjoint families.

For the following calculation we use the non-self-dual parameters
$(a,b,c,d,e,f)$ of that presentation. A useful further elimination illustrates the arithmetic structure. In the
non-self-dual case set $u=c-d$ and $P(u)=u^2-au-1$. Subtraction of the
first two equations and use of the third give
\[
 b(e-f)=P(u),\qquad dP(u)=-b^2u.
\]
If $b\ne0$, then $P(u)\ne0$: an integer root of $u^2-au-1$
would have $u=\pm1$, and $dP(u)=-b^2u$ would force $b=0$.
Since $P(u)\equiv-1\pmod u$, one has $\gcd(P(u),u)=1$; consequently
\[
 P(u)\mid b^2\qquad(b\ne0).
\]
The branch $b=0$ is retained separately rather than discarded by division.
This divisor restriction can replace an unstructured search by arithmetic
branches; it is not, by itself, a bound on the parameters.

\subsection{What remains to be proved}
A parametric classification of all rank-four fusion rings is a concrete
Diophantine problem with a fixed small system of equations, rather than
an unformulated search. The proper-subring families and the dimension-field
reductions proved here make a systematic approach plausible. Nevertheless,
the equations have unbounded integral solutions, and the present arguments
do not prove a descent to finitely many seed families in the residual
self-dual case.

For the Grothendieck-ring problem the more focused target is the residual
cubic or quartic branch of Question~\ref{question:unrestricted-simple}.
A uniform bound on its potentially categorifiable parameters, or a uniform
moment or arithmetic obstruction for every additional family, would close
the pivotal classification. General categorification requires, in addition,
control of the possible proper signed sphericalizations. The finite
certificate records the outcome of that control over the bounded
census; its uniform completeness remains a separate question.

Thus success of the fixed criterion suite at every tested multiplicity is
meaningful evidence for further investigation, but neither a proof that
all necessary criteria have been found nor a reason to identify raw fusion
ring enumeration as the sole remaining obstacle. The companion results
supply reductions toward both tasks while keeping their logical roles
separate.

The finite paper formulates the unrestricted fifteen-ring classification
and the stronger sufficiency of its certificate criteria as conjectures.
The reductions proved here support those conjectures but leave
Question~\ref{question:unrestricted-simple} open.
